On mélange initialement à $t_{0}=0$, une solution aqueuse de sulfate de fer(II)
(\ce{Fe$^{2+}$_{(aq)} + SO$_4^{2-}$_{(aq)}}) avec une solution aqueuse de
sulfate de mercure(II) (\ce{Hg$^{2+}$_{(aq)} + SO$_4^{2-}$_{(aq)}}).  Le volume
du mélange est $V=\qty{100}{\mL}$.

La transformation chimique qui se produit est modélisée par l'équation chimique
suivante~:
\begin{center}
    \ce{2Fe$^{2+}$_{(aq)} + 2Hg$^{2+}$_{(aq)} -> 2Fe$^{3+}$_{(aq)} + Hg$_2^{2+}$_{(aq)}}
\end{center}

\begin{donnees}
    Concentrations molaires effectives des ions dans le mélange à l'état
    initial~: \\
    $\left[\ce{Fe$^{2+}$_{(aq)}}\right]_{0}=\qty{1.0e-2}{\mol\per\L}$~;
    $\left[\ce{Hg$^{2+}$_{(aq)}}\right]_{0}=\qty{2.0e-2}{\mol\per\L}$.
\end{donnees}
\begin{minipage}{0.52\linewidth}
    Le suivi de la transformation permet de tracer la courbe représentant
    l'évolution temporelle de l'avancement $x$ de la réaction (figure
    ci-contre).

    \begin{enumerate}
        \item Dresser le tableau d'avancement de la réaction.
              \textbf{(1pt)}
        \item Vérifier que la valeur de l'avancement maximal est
              $x_{\max}=\qty{5e-4}{\mol}$.
              \textbf{(0,5pt)}
        \item Définir le temps de demi-réaction $t_{1/2}$ et déterminer sa valeur
              graphiquement.
              \textbf{(1pt)}
        \item Déterminer, en ($\unit{\mol\per\L\per\h}$), la vitesse volumique de
              réaction à $t_{0}=0$.
              \textbf{(1pt)}
        \item Comment évolue la vitesse volumique de réaction au cours de cette
              transformation~?
              \textbf{(0,5pt)}
    \end{enumerate}
\end{minipage}
\hfill
\begin{minipage}{0.44\linewidth}
\begin{tikzpicture}[]
    \def\ta{20}
    \def\xf{0.5}
    %\pgfmathparse{ln(2)*\ta}\let\tm=\pgfmathresult
    \pgfmathparse{\xf/\ta}\let\p=\pgfmathresult
    \begin{axis}[
            cycle list name=my color list,
            width=9cm,height=8cm,
            xlabel=$t~(\unit{\h})$,ylabel=$x~(\unit{\mmol})$,
            xmin=0,xmax=59,
            ymin=0,ymax=0.54,
            %xtick={},
            %xticklabels={},
            %ytick={},
            %yticklabels={},
            variable=t,
            samples=200,
        ]
        \addplot[domain=0:60] {\xf*(1-exp(-\t/\ta))};
        \addplot[domain=0:21,dashed,thick] {\p*\t};
    \end{axis}
\end{tikzpicture}
\end{minipage}



